%%%%%%%%%%%%%%%%%%%%%%%%%%%%%%%%%%%%%%%%%%%%%%%%%%%%%%%%%%%%%%%%%%%%%%%%%%%%%%%%
% DisasterSociety — IEEE journal working manuscript
% Rebuilt in the dedicated IEEEtran workspace on 2026-08-14.
% Quantitative findings are intentionally deferred to a separate evidence-sync pass.
%%%%%%%%%%%%%%%%%%%%%%%%%%%%%%%%%%%%%%%%%%%%%%%%%%%%%%%%%%%%%%%%%%%%%%%%%%%%%%%%
\documentclass[journal]{IEEEtran}

\usepackage{amsmath,amssymb}
\usepackage{graphicx}
\usepackage{multirow}
\usepackage{booktabs}
\usepackage{tabularx}
\usepackage{array}
\usepackage{placeins}
\newcolumntype{Y}{>{\raggedright\arraybackslash}X}
\usepackage{threeparttable}
\usepackage[hidelinks]{hyperref}
\usepackage[nameinlink,noabbrev]{cleveref}


\title{DisasterSociety: A Generative-Agent Platform for Simulating Resident and Collective Responses in Urban Disasters}

% PENDING(final submission metadata): replace with real authors and affiliations.
\author{Anonymous Author(s)%
\thanks{Author names, affiliations, funding information, and the corresponding
author will be added after the target IEEE journal and review policy are fixed.}}

\markboth{IEEE Transactions Journal Draft, August 2026}%
{Anonymous Author(s): DisasterSociety}

\begin{document}
\maketitle


%%%%%%%%%%%%%%%%%%%%%%%%%%%%%%%%%%%%%%%%%%%%%%%%%%%%%%%%%%%%%%%%%%%%%%%%%%%%%%%%
\begin{abstract}
Urban disaster response emerges from a sequence of socially situated decisions:
residents receive unequal information, interpret changing conditions, negotiate
with household members, revise plans after disruption, and act under shared
physical and resource constraints. A one-shot evacuation label cannot represent
these coupled information, cognition, coordination, and execution processes. We
present \emph{DisasterSociety}, a generative-agent platform for studying resident,
household, network, and collective response in a dynamic urban-disaster world.
The platform separates three forms of state that are often conflated: private
resident cognition, household-shared constraints and accepted commitments, and
authoritative physical state. A resident language model can propose a structured
intention, but only deterministic household and world logic can allocate
resources or execute it. Message disposition is likewise kept separate from
validated household commitment and physical execution. An \emph{EventPack}
interface supplies time-indexed hazard, warning, road, resource, and provenance
data, while an event-driven kernel connects limited observation, memory,
feedback, planning, interaction, batched arbitration, and execution in a closed
multi-step loop. We further define a multi-scale, multi-evidence evaluation
methodology that distinguishes software and process verification, empirical and
behavioral reasonableness, mechanism-specific ablation, and robustness and
efficiency evidence. A Carr-informed controlled case operationalizes the
framework without claiming historical reconstruction or population-level
prediction. The resulting research prototype is intended as an inspectable
computational framework for asking how heterogeneous resident actions become
household and collective disaster-response processes.
\end{abstract}

\begin{IEEEkeywords}
Generative agents, multi-agent simulation, urban disaster response, household
coordination, social simulation, event-driven simulation, mechanism ablation.
\end{IEEEkeywords}


%%%%%%%%%%%%%%%%%%%%%%%%%%%%%%%%%%%%%%%%%%%%%%%%%%%%%%%%%%%%%%%%%%%%%%%%%%%%%%%%
\section{Introduction}

Urban disaster response is not an isolated choice made at the instant a warning
arrives. Protective-action research distinguishes message receipt and
comprehension from threat appraisal, action appraisal, and the practical
impediments that shape whether and when a resident acts
\cite{lindell2012padm,mileti1992warnings,huang2012ike}. Wildfire studies likewise
show that warning exposure, preparation, and departure unfold over time rather
than collapsing into a single binary outcome
\cite{kuligowski2021wildfire,grajdura2021awareness}. The process is also social.
Household members may begin in different locations, coordinate gathering and
shared-vehicle use, decide whether to reunite, and depart together or in separate
groups \cite{liu2014uniting,maghelal2017split}.
Consider the controlled running example in \cref{fig:running-example}: Resident A
receives an official warning while Resident B does not; the two exchange
proposals about driver, companions, route, and departure time; after the primary
route closes, the original plan fails and the household must form a feasible
revision before the world permits execution. A useful simulation must preserve
every distinction in this chain---who observed what, which proposal was merely
sent, what the household actually accepted, and which action the physical world
ultimately executed.

\begin{figure*}[t]
  \centering
  \includegraphics[width=\textwidth]{../figures/fig1_running_example.png}
  \caption{A disaster evacuation is a coupled information, household, and
  execution process rather than a one-shot label. (A) Unequal warning exposure
  produces different private resident beliefs even within one household.
  (B) Messages can negotiate a route, vehicle, companions, and departure time,
  but only an accepted and feasible proposal becomes shared coordination state.
  (C) A road closure invalidates the original intention; feedback supports a
  revised plan, and world arbitration determines whether the revision executes.
  The scenario is Carr-informed and controlled; it illustrates the modeled
  mechanisms rather than reconstructing the historical Carr Fire timeline.}
  \label{fig:running-example}
\end{figure*}

Agent-based models provide a natural bottom-up representation of heterogeneous
actors and their interactions, while large language models (LLMs) make it
possible to express context-sensitive judgments and communication without
enumerating every behavioral rule. Generative Agents connected memory,
reflection, planning, and interaction in a persistent artificial community
\cite{park2023generativeagents}; subsequent systems have scaled societal and
urban simulation with LLM agents, in some cases through hybrid
LLM/numerical-agent designs,
\cite{zhang2025agentsociety,bougie2025citysim,zhou2026gasim}. Disaster-oriented
work now spans theory-informed individual evacuation prediction
\cite{chen2025flare}, LLM residents embedded in spatial fire or crowd simulators
\cite{dang2025fireagents,yang2026whenagents}, and integrated hazard--population
demonstrations \cite{sultimov2026respond}. These advances establish important
components, but they also expose three separations that a resident-society
simulation must make explicit. A dynamic world is not the same object as a
resident's belief about it; a household is neither an indivisible actor nor a
shared mind; and a linguistically plausible action is not necessarily an
executable action. The cited systems expose different subsets of these concerns
and do so at different levels of abstraction. Our goal is therefore not to
claim that any individual component is unprecedented, but to make their
interfaces explicit in one executable research framework.

We study the following problem. Given a time-indexed disaster scenario,
heterogeneous residents grouped into resource- and care-constrained households,
and an evolving physical environment, generate a sequence of resident
observations, private state updates, communications, intentions, and executed
outcomes from which individual, household, network, and collective processes can
be inspected. The problem has three hard requirements. First, information must
be exposure-bounded: co-residence cannot automatically synchronize inboxes,
memories, route beliefs, or unexpressed plans. Second, coordination must be
commitment-bounded: sending, delivering, or accepting words does not itself
reserve a vehicle, satisfy a care dependency, or make a departure time feasible.
Third, action must be world-bounded: all contemporaneous intentions must be
resolved against one authoritative snapshot so that closed roads, vehicle
capacity, member position, and competing resource claims have consistent
consequences. These requirements turn disaster response from a prediction task
into a stateful social-process simulation problem.

Satisfying the three requirements creates corresponding methodological
challenges. The first is to represent resident cognition across time without
leaking global state into every prompt: memory must contain processed
experience, feedback must reflect observed or executed outcomes, and planning
must remain an intention until execution. The second is to represent household
coordination without inventing a household-level language model: private
residents communicate through a lossy channel, while a separate household object
holds membership, resources, care obligations, and versioned commitments. The
third is evidential. Verification of a simulation implementation, empirical
reasonableness of its behavior, causal evidence from a mechanism ablation, and
robustness across stochastic or model conditions answer different questions
\cite{sargent2013verification,klugl2008validation,grimm2005pattern}. A credible
evaluation must preserve those distinctions instead of treating a unit test, a
single aggregate fit, or an LLM plausibility judgment as interchangeable proof.

DisasterSociety addresses these challenges through the architecture developed in
\cref{sec:method}. An \emph{EventPack} and scenario adapter organize external
hazard, warning, population, road, resource, and provenance inputs. The
event-driven kernel exposes each resident only to delivery- and
location-appropriate information; resident adapters maintain private memory,
feedback, and plans; the interaction layer advances messages through explicit
life-cycle states; the household layer validates and versions coordination
commitments; and the world batches all submitted intentions before committing
physical change. Crucially, the language model is used inside the resident
cognition boundary, not as the source of physical truth or the resolver of shared
resources. The same boundaries expose mechanism switches, state snapshots,
message and commitment histories, execution reasons, model-call metadata, and
terminal run status for evaluation and audit.

This paper makes three contributions. First, it presents DisasterSociety, a
configurable generative-agent platform that connects heterogeneous residents and
households to a time-indexed disaster world while preserving explicit
information, state, and execution boundaries (\cref{sec:method}). Second, it
defines a resident social-process model in which private cognition, lossy
interaction, household-shared constraints, versioned commitments, batched world
arbitration, and execution feedback form a closed multi-step loop
(\cref{sec:dynamic-loop}). Third, it develops a multi-scale, multi-evidence
evaluation methodology that maps claims about residents, households and
networks, collective outcomes, and system operation to appropriate system,
empirical, mechanism, and robustness evidence (\cref{sec:evaluation-method}). A
Carr-informed controlled case instantiates these contributions as an empirical
test bed; it is not presented as a historical reconstruction, an automatic city
generator, or a population-level evacuation predictor.


%%%%%%%%%%%%%%%%%%%%%%%%%%%%%%%%%%%%%%%%%%%%%%%%%%%%%%%%%%%%%%%%%%%%%%%%%%%%%%%%
\section{Related Work}

\subsection{Protective Action, Household Coordination, and Warning Diffusion}

Protective-action research treats disaster response as a staged process. In the
Protective Action Decision Model, environmental and social cues and warning
messages first have to be received, attended to, and understood; subsequent
threat and action appraisals remain subject to situational facilitators and
impediments \cite{lindell2012padm}. Studies of disaster warnings and hurricane
evacuation similarly distinguish communication from action and the decision to
evacuate from the timing of departure
\cite{mileti1992warnings,huang2012ike,kuligowski2021wildfire}. This literature
motivates our separation of a world event, a delivered observation, a processed
resident memory, a plan, and an executed outcome. It does not, by itself, dictate
the software state machine or establish that an LLM reproduces those processes.

Households add another level of organization. Multi-adult households may need to
gather members distributed across locations and allocate shared vehicles;
observed and stated plans do not always place all members in one group or at one
departure time
\cite{liu2014uniting,maghelal2017split}. At the inter-household level, agent-based
studies have represented networked information through different abstractions.
Nagarajan et al. simulate rule-governed neighbor warning dissemination, while
Roy et al. model how multiple information sources shape risk-perception dynamics
and evacuation decisions \cite{dawson2012warning,raihan2022information}. Our
focus is complementary: decision-capable residents retain separate cognition,
whereas the household is a shared constraint and commitment object, and the
message state is kept separate from resident memory and physical execution.

Evacuation simulation also has a substantial transportation and operations
lineage. Dynamic traffic models represent participation, departure time,
destination, route choice, network assignment, and changing capacity at
different levels of behavioral detail \cite{pel2012travel,murraytuite2013overview}.
This literature is indispensable for estimating clearance, congestion, and
traffic-management consequences, and it already recognizes time-varying
information and en-route adaptation. DisasterSociety does not replace a
high-fidelity traffic simulator. Its contribution is upstream and social: it
preserves the resident information and household coordination process that
produces departure parties and route intentions, after which a detailed network
model could implement the same arbitration interface.

\subsection{Generative Agents and Urban Social Simulation}

Generative-agent research provides mechanisms for persistent, language-mediated
behavior. Generative Agents combined an episodic memory stream, retrieval,
reflection, planning, and communication, and evaluated both perceived
believability and emergent social behavior with component ablations
\cite{park2023generativeagents}. Social Simulacra used populated prototypes to
support social-computing design \cite{park2022socialsimulacra}, while SOTOPIA
evaluates social intelligence through interactive goals and relationships
\cite{zhou2024sotopia}. These systems make interaction and persistent state
first-class objects, but their evidence concerns artificial communities, social
competence, or design prototyping; it cannot be transferred directly into a
claim about disaster behavior.

Recent platforms broaden the scale and environment of LLM-agent simulation.
AgentSociety integrates parallelized LLM agents with realistic societal
environments and interactions \cite{zhang2025agentsociety}; CitySim models
long-horizon urban schedules, goals, spatial memory, mobility, and aggregate city
patterns \cite{bougie2025citysim}; GASim combines LLM-driven core agents with
graph-accelerated numerical agents to accelerate memory retrieval and numerical
agent updates while reducing token use
\cite{zhou2026gasim}. Their contributions center on scalable societal or urban
simulation. DisasterSociety instead centers on short-horizon, disruption-driven
social processes where unequal information, household dependencies, resource
conflicts, and rejected actions must be represented without allowing resident
language models to overwrite physical truth.

\subsection{LLM Agents for Disaster and Evacuation}

LLMs have recently been introduced at several points in the disaster-modeling
pipeline. FLARE uses behavioral theory and LLM reasoning for individual wildfire
evacuation prediction on post-event surveys \cite{chen2025flare}. Its unit and
estimand are therefore different from a multi-step resident society, even though
the behavioral factors are relevant. Dang et al. embed LLM agents with memory,
cognition, communication, and movement in a cellular-automata fire environment
\cite{dang2025fireagents}. Yang et al. couple LLM decision histories and agent
communication with pedestrian and vehicle evacuation in a dynamic ABM
\cite{yang2026whenagents}. RESPOND joins flood forecasting, alerts, road
closures, social signals, and population agents in an integrated demonstration
\cite{sultimov2026respond}; generative-agent work has also addressed
socio-technical disaster recovery \cite{blagojevic2026sociotechnical}. These
studies show that individual prediction, spatial evacuation, dynamic hazards,
communication, and integrated platforms are all active research directions.

Our distinction is consequently an interface claim rather than a priority claim.
Across the systems compared above, resident-private cognition,
household-shared resources and care obligations, explicit message state,
batched world arbitration, and mechanism-linked evaluation are distributed
across different task definitions and abstraction levels. DisasterSociety makes
their joint interfaces explicit to study how resident actions become social
processes. It does not claim the first use of agents, households, warnings, LLM
memory, or evacuation simulation.

\subsection{Synthetic Household Populations}

Synthetic reconstruction connects aggregate small-area constraints to
microdata. Classical work uses iterative proportional fitting (IPF) to estimate
weights over household microdata and then selects donor households
\cite{beckman1996synthetic}. Carrying a whole donor household preserves observed
member relationships and attributes inside that record, which is useful when
shared vehicles, co-residence, and care dependencies must be instantiated
together. It does not imply that all person-level marginals are fitted.
Household--person methods such as iterative proportional updating and related
extensions explicitly introduce person controls while adjusting household
weights \cite{ye2009ipu,pritchard2012population,guo2007population}.
DisasterSociety's current population path fits selected household-level tract
margins from weighted donor records and reports person-level distributions as
diagnostics. We retain this distinction because a plausible household structure,
a calibrated person population, and validated resident behavior are separate
claims.

\subsection{Verification and Evaluation of Agent-Based Social Simulation}

Simulation credibility is purpose-dependent. Conceptual-model validity,
implementation verification, data validity, and operational validity address
different questions \cite{sargent2013verification}. ABM-specific methodology
accordingly combines, rather than conflates, face, calibration, sensitivity, and
statistical evidence \cite{klugl2008validation}. Pattern-oriented modeling
further argues that bottom-up mechanisms should be confronted with multiple
observed patterns at different levels instead of a single aggregate target
\cite{grimm2005pattern}. For stochastic models, systematic experimental design
and repeated runs are needed to characterize run-to-run variation
\cite{lorscheid2012blackbox}; local mechanism perturbations and global
sensitivity analyses answer different questions in nonlinear ABMs
\cite{tenbroeke2016sensitivity}.

Generative social simulation adds language-model and prompt dependence to these
longstanding concerns. Generative Agents uses human believability judgments and
component ablations \cite{park2023generativeagents}, while SOTOPIA provides an
interactive social-intelligence benchmark \cite{zhou2024sotopia}. Broader
social-science work identifies both opportunities and bias, replication, and
validity risks for generative AI \cite{bail2024generativeai}. These evidence
forms address different constructs; none alone verifies event causality,
household feasibility, or model robustness. We
therefore distinguish two orthogonal dimensions. \emph{Analysis scale} denotes
individual, household and network, collective, or system-level outcomes;
\emph{evidence type} denotes software and process verification, empirical and
behavioral reasonableness, mechanism ablation, or robustness and efficiency.
This organization is our evaluation design, not a claim that every cell is
already validated or that calibrated targets constitute independent evidence.


%%%%%%%%%%%%%%%%%%%%%%%%%%%%%%%%%%%%%%%%%%%%%%%%%%%%%%%%%%%%%%%%%%%%%%%%%%%%%%%%
\section{Method}
\label{sec:method}

DisasterSociety is a discrete-time, event-driven platform for studying how
resident-level cognition and communication produce household and collective
outcomes under a changing disaster environment. The method is organized around
contracts rather than a monolithic ``LLM simulator.'' A resident model may
interpret information and return a typed intention but cannot directly mutate
shared resources or physical state. Interaction and household layers validate
responses and maintain versioned coordination records under adapter-specific
checks. Deterministic household and world code then allocate shared resources
and apply execution state; a stored commitment is therefore a validated
coordination record, not a guarantee of global feasibility before world
arbitration. This decomposition is documented consistently with ODD's emphasis
on explicit entities, state variables, process scheduling, observation, and
stochasticity \cite{grimm2020odd}. It also lets the
same event-driven kernel host two Carr-informed adapters: an empirical-process
path with richer resident profiles and departure parties, and a controlled
mechanism path with switchable memory, feedback, planning, and interaction.
They share causal and state boundaries but do not falsely share one internal
cognitive representation.

\begin{figure*}[t]
  \centering
  \includegraphics[width=\textwidth]{../figures/fig2_system_architecture.png}
  \caption{Static architecture of DisasterSociety. Time-indexed external events,
  household donor records, and controlled case configuration enter an
  event-driven engine. Resident-private cognition, the explicit interaction
  channel, household-shared constraints and commitments, and world physical
  facts are separate state owners. A resident produces a structured intention;
  deterministic household and world logic produces the execution outcome.
  Per-step traces preserve world, resident, interaction, intention, and outcome
  records; the controlled adapter additionally records a configuration and
  selected-source fingerprint. The diagram shows
  the shared platform contracts; adapter-specific cognition and coordination
  implementations are detailed in the text.}
  \label{fig:architecture}
\end{figure*}

\subsection{Problem Formulation and State Boundaries}

Let residents be indexed by $i\in\mathcal{R}$ and households by
$h\in\mathcal{H}$, with a membership function $g(i)=h$. At simulation step $t$,
the state is
\begin{equation}
  S_t=\left(W_t,\{C_{i,t}\}_{i\in\mathcal{R}},
  \{H_{h,t}\}_{h\in\mathcal{H}},Q_t\right),
  \label{eq:state}
\end{equation}
where $W_t$ is authoritative world state, $C_{i,t}$ is resident $i$'s private
cognitive state, $H_{h,t}$ contains shared household constraints and accepted
coordination state, and $Q_t$ contains events and adapter-specific message,
receipt, and ledger state. These are separately owned namespaces. Updating
$C_{i,t}$ does not imply an update to $C_{j,t}$, and neither a private plan nor a
private belief becomes household or world state without an explicit protocol
transition. Thus, two co-resident adults may hold different warnings or route
beliefs; an unexpressed private plan is not a household agreement; and a
resident's belief that a road is open cannot reopen it in the physical world.

An access function $\Omega_i$ maps world state and delivery records to the
resident's bounded observation,
\begin{equation}
  O_{i,t}=\Omega_i(W_t,Q_{\leq t};A_{i,t}),
\end{equation}
where $A_{i,t}$ encodes channel, recipient, location, and adapter-specific
access rules. Resident $i$ returns a typed intention
\begin{equation}
  I_{i,t}=\Pi_i(C_{i,t},O_{i,t};\theta,\phi),
  \label{eq:intention}
\end{equation}
where $\theta$ denotes the declared model backend and $\phi$ the enabled
mechanisms and prompt protocol. $I_{i,t}$ is a proposal, not a state mutation.
After decisions, an interaction transition processes the current inbox and
returned message responses into intermediate private and coordination state,
\begin{equation}
  (\widetilde C_t,\widetilde H_t,\widetilde Q_t)=
  \mathcal{P}(C_t,H_t,Q_t,\{I_{i,t}\}).
  \label{eq:interaction-transition}
\end{equation}
The world then resolves all contemporaneous intentions against one physical
snapshot and the intermediate coordination state,
\begin{equation}
  \{R_{i,t}\}=\mathcal{A}\!\left(W_t,\{\widetilde H_{h,t}\},
  \{I_{i,t}\}\right),
  \label{eq:arbitration}
\end{equation}
and returns a keyed execution outcome $R_{i,t}$ containing status, executed
action, rejection reason, resource allocations, and state delta. Only keyed
outcomes can commit execution-derived resident fields or physical movement;
other cognitive and coordination transitions occur through delivery,
observation, and interaction processing.

The four principal objects in \cref{fig:architecture} therefore have distinct
ownership. A \emph{Resident} is a cognitive and decision unit. A
\emph{Household} is a system object that owns membership, vehicles, care
requirements, and accepted commitments; it is not a household-level LLM. The
\emph{Interaction} layer owns message state. The \emph{World} owns authoritative
hazard, route, location, and movement arbitration, while the Household owns
membership, commitments, care, and household-resource state. Neither is
writable by the resident model. Resident attributes also distinguish
\texttt{decision\_capable}, a persistent eligibility property;
\texttt{decision\_policy}, which may be generative, rule-based, or dependent;
and a step-specific wake gate. A dependent member can retain location, safety,
caregiver, and transit state without making a language-model call.

\subsection{Scenario Data and the EventPack Interface}

An \emph{EventPack} is a provenance-aware scenario data package rather than an
LLM prompt or an assertion that every experiment uses the same raw file format.
Its validated manifest declares an event identifier, hazard type, time zone,
spatial reference system, simulation step length, named data assets, and a
provenance record. The current loader validates manifest fields and the
existence of declared assets and a declared provenance file; it does not yet
verify content checksums or the internal provenance schema. A scenario adapter
converts those assets, or a separately frozen experimental configuration, into
the small
\texttt{EventSource} contract consumed by the kernel. This contract returns the
exogenous events due at step $t$; warnings, orders, hazard updates, and road
events are not generated by resident language models.

The separation between data package and adapter is important in the current
prototype. The controlled mechanism path loads its Carr EventPack through the
validated interface and references a manifest-declared route asset, although the
current runner resolves that configured path directly rather than enforcing
manifest membership through the asset accessor. The empirical-process path
assembles official-order and route-closure events from frozen profiles,
calibration quantities, and scenario
configuration before passing them to the same kernel contract. Consequently,
EventPack is the reusable platform interface, but we do not claim that all
existing runners already pass through one identical typed loader.

The Carr-informed case is deliberately a controlled scenario. Its final fire
perimeter is a single endpoint, not a historical hazard time series; its two
routes are derived experimental assets, and the road disruption is a controlled
perturbation rather than an observed Carr closure. The world may therefore hold
hazard distance, route state and capacity, resident and dependent locations,
vehicle state, care obligations, and household commitments, while the paper
does not describe the case as a reconstruction of the historical fire or of
historical traffic operations. Neither current adapter converts the final
perimeter into a time-varying fire surface: both use declared scalar
hazard-distance trajectories and controlled route closures, with the sourced
perimeter retained as case provenance.

\subsection{Household Population and Resident Initialization}

The population pipeline starts from weighted household microdata. Public-use
donor weights initialize the household distribution; iterative proportional
fitting then adjusts donor weights within tract for four groups of household
margins: income, vehicle availability, household structure, and older-adult
presence. Sampling is atomic at the donor-household level: when a donor is
selected, its member records are copied together and receive new synthetic
identifiers. This retains internally observed relationships and compatible
household attributes, following the established household-IPF tradition
\cite{beckman1996synthetic}. It does not imply person-level calibration.
Age, education, employment, and other member distributions are therefore
reported as diagnostics unless explicitly fitted. Full household--person
procedures such as IPU address a stronger target and are not claimed as the
current implementation \cite{ye2009ipu,pritchard2012population,guo2007population}.

Initialization converts each sampled household into member, resource, and care
state. In the current Carr initialization, age 18 or older defines the
transparent decision-capable baseline. Functional limitations and need for
execution assistance are separate care properties and do not automatically
remove an adult's decision capability; minors are dependent execution members.
E1 prefers the public-use reference person as coordinator, with a stable
adult-ID fallback, whereas E2 assigns the first ordering-key adult as the
controlled coordinator. Vehicles have household ownership, location, capacity,
occupancy, and reservation state. E1 derives vehicle count from the profile but
assigns a common five-seat capacity. E2 selects exactly-one-vehicle households
and sets capacity at or above household size, so its mechanism experiment does
not estimate empirical vehicle-capacity scarcity. The richer empirical-process
adapter combines these household records with static and latent profile fields
sampled only from the designated Carr training data without using evacuation
outcomes; it records field-level source labels so donor attributes and survey-
informed traits remain distinguishable.

\subsection{Event-Driven Kernel and Causal Step Order}
\label{sec:dynamic-loop}

The kernel executes ten per-step operations, followed by run-level terminal
classification (\cref{fig:dynamic-loop}). It (1) advances the clock and world,
(2) applies due exogenous events, (3) creates official warning or order
receipts, (4) evaluates the wake gate, and (5) collects decisions concurrently
while physical world state is immutable. It then (6) processes the current
inbox and returned message responses, potentially updating private memory and
household coordination state; (7) wraps all decisions, resolves the complete
physical-action batch against one world snapshot, and applies the resolved
actions; (8) commits keyed execution outcomes and emits outgoing messages; (9)
advances message delivery; and (10) records the step snapshot. After the loop or
an exception, the run is classified as \textsc{valid}, \textsc{invalid}, or
\textsc{aborted} under the frozen operating protocol.

\begin{figure*}[t]
  \centering
  \includegraphics[width=\textwidth]{../figures/fig3_dynamic_loop.png}
  \caption{Dynamic method loop and its two explicit state chains. The upper
  loop separates limited observation, private memory, belief and feedback
  update, planning, household interaction, and world arbitration. The message
  chain illustrates the controlled E2 recipient-receipt protocol by
  distinguishing sending, delivery, processing, response, and validated
  commitment formation. The action chain distinguishes an intention from
  batched resolution, executed or rejected outcome, and subsequent state
  update. Mechanism toggles in the figure denote the controlled E2
  operationalization; the empirical E1 adapter shares the kernel and state
  boundaries but uses the private-state representation described below.}
  \label{fig:dynamic-loop}
\end{figure*}

The order prevents several common causal errors. A resident cannot act on a raw
warning it did not receive, because delivery precedes waking and the prompt
reads the private inbox rather than the event list. All decisions observe the
same physical phase, and typed outcomes precede execution-derived resident
updates. The controlled E2 world further verifies that rejected movement does
not consume or mutate the corresponding physical resource. The current E1 world
uses the same typed boundary but still has audit findings around unknown-route
handling, capacity bookkeeping after party rejection, and out-of-party intent
matching; until those checks are corrected, this invariant is claimed for the
verified contracts and E2 path, not for every malformed E1 input.

Message timing and receipt granularity are adapter-specific within this common
schedule. In the controlled interaction implementation, household direct
messages and community messages have separate delivery models,
recipient-specific receipts, and future processing phases. In the empirical
adapter, one social-message record may address several recipients and stores a
shared status and timestamps, although delivered copies enter particular
resident inboxes. Both message classes advance at step end and are processed
when a recipient next wakes. For messages with recipients, enqueueing alone
does not update recipient memory. E1 additionally has a zero-recipient
self-acceptance path and may form a solo party from an evacuation intention;
neither path is evidence of message-mediated household coordination.

\subsection{Two Carr-Informed Resident Adapters}

The shared kernel deliberately supports different operationalizations of
resident cognition. Treating the two current adapters as one implementation
would obscure which claims each can support.

\paragraph{Empirical-process adapter (E1).}
Each decision-capable adult has a private inbox, a bounded recent-memory list,
perceived routes, a private view of commitments, and recent execution feedback.
The prompt exposes a richer resident and household profile, including household
members, functional limitations, available vehicles, housing, and relevant
care context. Official orders are assigned according to frozen empirical quotas
and timing profiles and then delivered directly to preassigned recipients; they
are not redrawn from the controlled adapter's Bernoulli delivery model. When a
recipient processes an order or message, a processing timestamp and source
identifier are retained. Current route belief changes after a closed-route
execution rejection; the present implementation does not infer structured route
truth from arbitrary household or community text.

E1 represents a structured \texttt{DepartureParty} with traveler IDs,
accompanying dependents, caregiver assignments, vehicle, route, and departure
step. The current acceptance checker verifies household membership, acceptance
for traveler IDs marked decision-capable, caregiver inclusion, known vehicle,
vehicle capacity and location, a nonempty route, valid time fields, and active-
party overlap. Route existence and openness are intended to be checked by the
world. Execution requires matching intentions from every listed traveler ID. A
valid solo adult
intention may form a solo party without prior household negotiation; hence the
adapter studies coordination when it occurs but does not force every departure
through a household dialogue. E1 currently has no persistent executable
\texttt{current\_plan}; its private-process record instead contains recent
memories, commitments, and execution feedback.

The formal E1 configuration takes the first $n$ whole households from a frozen,
ordered profile file rather than resampling members independently for each run.
Hamilton largest-remainder allocation converts calibrated order-sequence
proportions into household counts; a named run-seed stream shuffles household
IDs before assigning voluntary-only, mandatory-only, both-order, or no-order
sequences. Orders are addressed to the household's designated focal resident.
In the current 30-minute-step configuration, voluntary and mandatory orders are
issued at steps 3 and 7, the controlled primary-route closure occurs at step 9,
and household hazard distance follows a linear approach from 9,000 m at 250 m
per step with a seeded household offset in $[-250,250]$ m. These quantities are
scenario controls, not reconstructed Carr timestamps.

E1 constructs an undirected decision-resident graph from three overlays:
household cliques, up to six nearest neighbors within 1,500 m, and a
Watts--Strogatz friend ring with degree four and rewiring probability 0.1.
Household messages use the household channel, whereas community messages are
routed to graph neighbors before delivery sampling. The run currently rebuilds
but does not persist the exact edge list, and missing geographic locations fall
back to a process-hash-derived coordinate; the graph is therefore a modeled
communication substrate rather than observed Carr social ties, and current E1
artifacts do not guarantee cross-process graph identity.

The E1 social-message ledger does not yet provide recipient-specific proof that
every addressed traveler accepted. One response can set the shared message
status to accepted before party feasibility is known, and the current
\texttt{accepted\_by} construction includes all addressed recipients. Message
disposition, party-acceptance result, and execution record must therefore be
audited separately; the current runner does not establish recipient-level
mutual acknowledgment or multi-adult message-mediated coordination.

An E1 resident wakes only when decision-capable and at least one of four
conditions holds: its inbox is nonempty, recent execution feedback is present,
a departure party involving it is due, or its household-specific hazard
distance is below 3,000 m. Consequently, a message can be delivered yet remain
unprocessed until a later wake opportunity; delivery counts and processed
counts have different denominators.

The current E1 implementation has an additional terminal-state defect that
prevents confirmatory use of its existing run artifacts. Recent execution
feedback is appended after each decision but is not consumed by the wake gate,
and the gate does not exclude a resident already in the safe zone. Once first
activated, a resident can therefore remain awake in later steps and continue to
send messages after its first safe arrival. The current party path also does not
recheck traveler origin or prior safe arrival, so repeated executed departures
by one already-safe resident are possible. Such records are invariant
violations, not distinct evacuations or split departures. E1 behavioral,
communication, departure, and efficiency quantities consequently require a
corrected terminal-state contract and a fresh run; the existing artifacts are
retained only for engineering diagnosis.

\paragraph{Controlled-mechanism adapter (E2).}
The controlled adapter gives paired conditions the same household cohort,
initial private primary-route plan, time-indexed disruption, and component-
isolated exogenous streams. Four code-level flags enable or remove memory,
feedback, planning, and interaction. The resident prompt contains its private
route observation, inbox, accepted household commitments, the same controlled
initial plan in every condition, and recent execution failures when feedback is
enabled. With planning enabled, a model decision may replace or revise the
initial plan; with planning disabled, it remains observable but cannot be
replaced through that update path. A true closure always affects world
arbitration. In feedback-enabled conditions, the next decision copies changed
world route state into the private view; rejected execution may additionally
provide bounded failure feedback. With feedback disabled, the pre-closure route
view persists. An \texttt{ExecutablePlan} records route,
vehicle, departure step, accompanying dependents, status, and last failure, and
plan history preserves revisions.

An E2 resident does not wake after reaching safety or when it is not a decision
unit. Otherwise it wakes for a nonempty inbox, a noncompleted plan, any accepted
household commitment, or hazard distance below 3,000 m. Because the controlled
initial plan is present in every paired condition, this gate provides a common
opportunity to decide while the planning flag changes only the plan-update
path. Community audiences in E2 expand to all controlled residents before
delivery sampling; this is a broadcast audience, not an explicit network
topology or a test of network diffusion.

Memory in the publication-facing controlled adapter has a deliberately narrow
operational meaning: delivery-gated warning retention and model-emitted memory
writes provide a cross-step historical signal. The platform also implements a
general bounded memory with rule-assigned importance, protected warning,
injury, and family items, low-importance/old-first eviction, and
importance--decay top-$k$ retrieval; these protection and eviction rules are
DisasterSociety-specific. The broader use of persistent ranked memory is
inspired by Generative Agents \cite{park2023generativeagents}. That full
retrieval path is used by the
base resident interface, but it is not presented as the exact mechanism tested
by the current controlled Carr adapter. This distinction keeps platform
capability and experiment-specific manipulation aligned.

\begin{table*}[t]
\centering
\footnotesize
\caption{Current Carr adapter contracts. Both adapters use the shared kernel,
typed intention--outcome boundary, and household/world state separation, but
their cohorts, private-state representations, communication records, and
estimands are intentionally different.}
\label{tab:adapter-contracts}
\begin{tabularx}{\textwidth}{@{}p{0.15\textwidth}Y Y@{}}
\toprule
Contract & E1 empirical-process adapter & E2 controlled-mechanism adapter \\
\midrule
Purpose and cohort & Carr-informed full-condition process study with richer
PUMS/Carr trait profiles; profile ordering and formal cohort are frozen &
Mechanism identification with seed-selected households containing exactly two
adults, at least one dependent, and exactly one reported vehicle \\
Hazard and roads & Household-offset scalar hazard trajectory, frozen order
timing, and controlled primary-route closure & Common scalar hazard trajectory,
controlled closure, and paired private primary-route initialization \\
Private cognition & Inbox, recent memory list, perceived routes, private
commitment view, and execution feedback; no persistent executable plan &
Mechanism flags, warning-history signal, feedback-gated route view, executable
plan and plan history, accepted commitments, and recent failure feedback \\
Communication & Recipient-specific official receipts; social messages use a
message-level record over a recipient list and a constructed social graph &
Recipient-specific receipts; household direct and community-audience delivery
are distinct, but the community audience is not a social-network topology \\
Coordination & Structured departure parties with care, vehicle, member, time,
and overlap checks, plus an allowed solo auto-party path & Versioned controlled
household proposals and commitments for the two decision members, with final
execution checks in the world \\
Vehicle setting & Profile-derived vehicle count with common five-seat capacity
in the formal configuration & One vehicle whose capacity is at least household
size; capacity scarcity is intentionally nonbinding \\
Trace and inference & Domain ledgers for orders, messages, commitments,
departures, resident state, and profiles; empirical reasonableness and process
constraints, subject to the current E1 ledger limitations & Per-step state and
call trace plus configuration/model/selected-source fingerprint; paired
full-versus-drop-one mechanism contrasts \\
\bottomrule
\end{tabularx}
\end{table*}

\subsection{Bounded Cognition and Structured Intentions}

All resident adapters obey an information boundary even though their JSON
payloads differ. A decision may include a profile, household facts visible to
that member, received inbox items, privately perceived hazard or routes,
selected private history, accepted commitments visible to the resident, and
recent execution outcomes. It never receives authority to edit the world. In
particular, co-residence does not synchronize inboxes, memories, beliefs, or
unexpressed plans, and the full world route snapshot is not copied into the E1
prompt.

The model returns a typed decision containing an assessment, an action from a
closed action set, optional destination, route, vehicle, and departure step,
accompanying dependents, outgoing messages, responses to received message IDs,
and optional memory writes. Schema validation constrains enums, field types,
and required structure; it does not prove that a syntactically valid identifier
exists. The controlled adapter deterministically rejects unknown route and
vehicle references. E1 validates household vehicle membership and a nonempty
route string, but does not yet offer a generic guarantee for destination IDs,
social-message recipient IDs, or unknown routes; an unknown E1 route can
currently abort instead of returning a clean rejection. These are implementation
limitations, not properties established by JSON validation. Malformed output
receives one schema-guided repair attempt before the gateway records a logical
failure.

\subsection{Messages and Versioned Household Commitments}

Messages carry sender, recipient or audience, channel, content, kind, structured
payload, sending step, and lifecycle metadata. At the platform-contract level,
implemented with recipient-specific receipts in the controlled adapter, the
guaranteed lifecycle prefix is
\begin{equation}
  \textit{sent}\rightarrow\textit{delivered}\rightarrow
  \textit{processed}.
  \label{eq:message-lifecycle}
\end{equation}
Only if the resident returns an explicit response may a processed receipt take
the disposition transition
\begin{equation}
  \textit{processed}\xrightarrow{\text{explicit response}}
  \{\textit{accepted},\textit{rejected}\};
  \label{eq:message-disposition}
\end{equation}
otherwise it may remain processed. The states are not rhetorical labels:
delivery determines inbox membership, and processing controls when the message
enters private history. E1 official information also uses recipient-specific
receipts, but E1 social messages currently use one message-level status and
timestamp record for a recipient list. Its social-message funnel must not be
interpreted as recipient-level reach or mutual acknowledgment, and an accepted
social-message status records response disposition rather than proof that party
acceptance succeeded. In particular,
$\textit{sent}\not\Rightarrow\textit{delivered}$ and
$\textit{accepted words}\not\Rightarrow\textit{executable action}$.

The controlled adapter validates proposal structure and timing on emission and,
on acceptance, verifies the two household decision members, a known household
vehicle, nonempty route and vehicle fields, and departure-time legality. Route
openness, movement compatibility, dependent care, and capacity are world-
arbitrated. E1 party acceptance adds member, caregiver, vehicle-capacity and
location, time, and overlap checks, but currently checks route syntax rather
than route membership or openness; execution supplies a separate check. When a
newer validated commitment revises an existing one, the prior version is
cancelled but retained
with a supersession link. The shared household commitment is authoritative;
private member views need not be automatically synchronized. In the controlled
adapter, the separate records support message, commitment, and execution
funnels. In E1, message-level acceptance is not treated as mutual acknowledgment
or feasible coordination without independent party-acceptance evidence.

\subsection{World Arbitration and Execution Feedback}

For each batch, the world converts intentions into resolved actions and checks
them against the physical snapshot. The Carr-controlled world checks route
openness and capacity, household vehicle state, dependent accompaniment, and a
compatible household commitment where required. Contending vehicle requests
are resolved using stable request attributes and a named random stream only as
the final tie-break, preventing condition labels from changing the exogenous
draw. The empirical world independently maintains true routes and capacities,
matches due departure parties to traveler intentions, and applies household
movement only after party checks succeed. Its batch processing does not claim
the identical seeded ordering policy used by the controlled adapter.

A keyed execution outcome is the only bridge from physical arbitration to
execution-derived resident fields and physical movement. Private cognition and
coordination may also change through official delivery, bounded route
observation, inbox processing, response handling, and memory updates. An
executed evacuation moves the verified residents and dependents and updates
vehicle and safety state. Commitment validation can reject a proposal for
member, caregiver, vehicle, time, or overlap reasons. World arbitration can
separately reject an executable intention for a closed route, route capacity,
vehicle conflict, future departure, missing traveler intention, or dependent-
accompaniment failure. When feedback is enabled, recent reasons and the current
plan's
\texttt{last\_failure} become bounded input to a later decision, allowing the
resident to revise its intention without granting the LLM direct world access.
The resulting loop resembles reasoning with action feedback
\cite{yao2023react,shinn2023reflexion}, but checked feasibility comes from
deterministic validators rather than self-correction language alone. The E1
arbitration limitations identified above remain outside this verified claim.

\subsection{Reproducibility, Failure Handling, and Trace}

All language-model calls pass through a gateway that controls concurrency,
timeouts, transient retry, budgets, schema validation, and one repair attempt.
The current cache key includes the declared model alias, serialized messages,
logical seed, and temperature; resolved provider identity and schema are not
part of that key and remain separate provenance concerns.
For providers with a narrower seed range, the logical seed is deterministically
mapped and both values are logged. Successful and cache-hit decisions record
model alias, seeds, status, token counts, estimated cost, and latency fields;
failure and fallback records preserve status and reason but do not populate
every success metric. A successful schema repair is not separately tagged in
the final record. Controlled E2 non-LLM randomness uses a component-isolated key
\begin{equation}
  k=(s_{\mathrm{run}},\,n_{\mathrm{stream}},\,i,\,t,\,d),
\end{equation}
where $d$ is a draw index; the experiment condition name is excluded so paired
conditions can share exogenous draws. E1 uses the same helper for several
components, but its current graph fallback and delivery counter still depend on
Python's process hash and therefore do not guarantee cross-process replay.

Failed model calls either abort under the frozen protocol or enter a logged
rule fallback whose intention remains subject to world arbitration. The
fallback rate is
\begin{equation}
  r_{\mathrm{fb}}=
  \frac{n_{\mathrm{fallback}}}
  {n_{\mathrm{ok}}+n_{\mathrm{cache}}+n_{\mathrm{fallback}}},
  \label{eq:fallback}
\end{equation}
while unrecovered failures are reported separately. A step record stores the
world snapshot, awake residents, typed intentions, keyed execution outcomes,
resident state, and interaction state; a terminal summary stores completion and
gateway statistics. The controlled runner additionally writes configuration,
selected-source, model-registry, and repository fingerprints. The empirical runner writes
official-receipt, message, commitment, departure, resident-state, and profile
ledgers, but the current prototype does not claim an identical provenance
bundle or cross-process byte identity for both adapters.

\subsection{Evaluation Methodology}
\label{sec:evaluation-method}

Evaluation is organized along two orthogonal dimensions. \emph{Analysis scale}
identifies whether a quantity concerns an individual, a household or network,
a collective population, or system operation. \emph{Evidence type} identifies
whether a result verifies implementation and causal contracts, assesses
empirical or behavioral reasonableness, estimates a mechanism-specific removal
effect, or measures robustness and efficiency. The distinction prevents a
deterministic replay from being presented as behavioral validity, a calibrated
warning margin from being reused as independent validation, or a large run from
being treated as proof of realism.

\begin{table*}[t]
\centering
\caption{Claim--evidence map used to design the evaluation. Scales and evidence
types are independent; the table gives the primary combination rather than a
requirement to populate every possible cell.}
\label{tab:evidence-map}
\begin{tabularx}{\textwidth}{@{}p{0.20\textwidth}Y Y@{}}
\toprule
Claim family & Primary quantities and analysis scale & Evidence required and
interpretive boundary \\
\midrule
Process and software correctness & Event order, delivery and processing
lifecycle, intention--execution invariants, resource conservation, terminal
status, and replay at system and household scales & Executed tests, invariant
checks, logs, provenance, failure and fallback accounting; does not establish
human behavioral validity \\
Empirical reasonableness & Warning/channel exposure, departure timing,
household constraints and group differences at individual, household, and
collective scales & Independent or prospectively held-out targets where
available; otherwise explicitly bounded respondent-sample references,
constraint violations, and planned human or expert judgment; calibrated inputs
are not independent confirmation \\
Mechanism function & Module-specific manipulation checks and downstream
household or collective consequences & Paired full-versus-removal runs with
predeclared direction and opportunity-to-act checks; estimates a conditional
removal effect, not every module interaction \\
Research utility and robustness & Seed, model, prompt, scale, latency, token,
cost, and failure variation at run and system scales & Multiple independent
runs and declared sensitivity conditions; operational scale alone is not
realism \\
\bottomrule
\end{tabularx}
\end{table*}

The Carr case uses two non-interchangeable populations. Carr-R retains a
version-frozen, de-identified respondent table for bounded empirical references
and auxiliary development diagnostics. Its former test partition has been
opened for diagnostic analysis and is not an untouched holdout. Carr-S contains
synthetic households for multi-step social-process
experiments; its agents are not one-to-one replicas of Carr-R respondents. A
survey field is assigned in advance to runtime input, calibration,
evaluation-only, or exclusion, and the same field cannot be both a run input and
its evaluation target in one experiment.

E1 uses train-designated, non-outcome-conditioned Carr trait blocks, but the
same survey also supplies some aggregate outcome and timing references. These
comparisons assess Carr-informed reasonableness rather than independent external
validation. Overall official reception is a calibration quantity, and
multi-select channel prevalence is not a first-heard distribution. The small
non-evacuator group and the absence of completed blinded human or expert review
further bound behavioral-validity claims.

For controlled mechanism evaluation, each module has a v8 runner-aligned check.
Memory uses conditional retention of a delivered official warning. Feedback
uses the share of residents with a pre-closure primary-route plan who later
record a post-closure alternate-route update. Planning uses the household share
that completes a post-closure alternate-route replan. Interaction uses the
household share that forms, after closure, an all-decision-member commitment
referencing the controlled alternate route, a known household vehicle, and an
integer departure step. These are scenario-specific manipulation checks rather
than general psychological-validity measures; the current runner does not
separately compute unsupported contradiction or update-delay metrics.

Downstream quantities include all-member safe location, closed-route rejection,
multi-adult co-departure, dependent safety or left-behind rejection, and shared-
resource conflict. For the forthcoming E1 results analysis, quantities will be
recomputed from initialized member identities, final authoritative locations,
recipient official-order receipts, and departure-party ledgers rather than
accepted from existing summary labels. Its present
\texttt{complete\_safe\_household\_departure} summary counts any household with
an executed party, while \texttt{coordinated\_household} may include a
one-decision-adult solo party; the current split and social-message acceptance
fields are likewise intermediate labels. None is eligible for Results until the
definitions below are implemented in a hashed, auditable analysis artifact. A
mechanism cannot receive credit from a metric it has no opportunity to affect.

The E1 analysis will keep analytically distinct denominators. Resident safe
arrival is the fraction of unique initialized decision-capable resident IDs
whose authoritative final location is the safe zone. Any-party household
departure requires at least one first-time valid arrival, whereas complete
household departure requires every initialized member to end in the safe zone.
Multi-adult coordinated departure is evaluated only among
households with at least two decision-capable adults and requires at least two
such adults to execute in the same party; a solo auto-party does not qualify.
Split departure requires at least two distinct household members to make their
first valid safe-zone arrivals in different executed parties or at different
departure steps. Multiple departure records for the same already-safe resident
are invariant violations, not split departure. Official-order reach is counted
from recipient receipts, while the current social-message funnel can be
reported only at message level until recipient receipts are implemented.
Message response, commitment feasibility, and physical execution rejection are
separate quantities. Order-to-departure delay is evaluated only for households
with both a delivered official-order receipt and a first-time valid household
departure, and is converted using the 30-minute step. Households with no
delivered order, no valid departure, or departure before the first delivered
order are reported as separate categories; negative lags are not truncated.
These definitions prevent current summary labels from silently changing the
estimand.

The controlled contrast for module $m$ and paired seed $s$ is
\begin{equation}
  \Delta_{m,s}=Y_{s,\mathrm{full}}-
  Y_{s,\mathrm{full}-m}.
  \label{eq:ablation}
\end{equation}
The estimand is the conditional effect of removing $m$ while the other modules
remain present; it is not a full factorial decomposition and does not identify
all interactions. The smallest run-level observation is a complete
$\text{seed}\times\text{condition}$ run, not a resident, household, or
resident-step. Full and removal cells sharing a seed form one paired contrast;
independent replication for that contrast is across distinct seeds, not between
the two cells within a seed. Differences are first formed within each paired
seed and only then summarized across seeds with uncertainty and sensitivity
analyses. This
evaluation design defines what evidence the experiments must supply; it does
not pre-announce their outcome.


% Experiments, Results, Discussion, and Conclusions are intentionally omitted
% from this pre-experiment working draft. They will be inserted only after a
% separate evidence-state synchronization and metric audit.


%%%%%%%%%%%%%%%%%%%%%%%%%%%%%%%%%%%%%%%%%%%%%%%%%%%%%%%%%%%%%%%%%%%%%%%%%%%%%%%%
\FloatBarrier
{\small
\bibliographystyle{IEEEtran}
\bibliography{references}
}

\end{document}
